%--------------------------------------------------------------------------------
%
%
%--------------------------------------------------------------------------------
\pagebreak
\unnumberedSection{IxTLB - Insert TLB Entry}

\begin{iPageDescEntry}{Syntax}
    \texttt{IITLB RegB, RegA} \\
    \texttt{IDTLB RegB, RegA}
\end{iPageDescEntry}

\begin{iPageDescEntry}{Format}
    The \texttt{IITLB} instruction uses the following instruction format. \\

    \begin{flushleft}
    \drawInstrFormat
        {0,1,3,5,6.5,8.5,9.5,11.5,16}
        {\OpCodeGrpSYS, \OpCodePTLB, 0, 0, RegB, 0, RegA, 0}
    \end{flushleft}
    
    The \texttt{IDTLB} instruction uses the following instruction format. \\

    \begin{flushleft}
    \drawInstrFormat
        {0, 1,3,5,6.5,8.5,9.5,11.5,16}
        {\OpCodeGrpSYS, \OpCodePTLB, 0, 1, RegB, 0, RegA, 0}
    \end{flushleft}

\end{iPageDescEntry}

\begin{iPageDescEntry}{Description}
    The \texttt{IxTLB} instruction inserts a new entry into the translation lookaside buffer (xTLB). The virtual address is stored in \texttt{RegB} and additional control information is provided by \texttt{RegA}. This is a privileged instruction.
\end{iPageDescEntry}

\begin{iPageDescEntry}{TLB Arguments}
	\texttt{RegB}  and \texttt{RegA} contain the arguments passed to the TLB subsystem. \texttt{RegB}  contains the virtual address. \texttt{RegA} contains the flags, access rights, page size and the physical address of the address range. 
	
	\begin{flushleft}
	\drawInstrFormatSixtyFour
        {0, 3, 13,16}
        {0, VPN, 0}
    \end{flushleft}
    \begin{flushleft}
	\drawInstrFormatSixtyFour
        {0,0.25,0.5, 0.75, 1, 1.25, 1.5, 1.75, 2, 3, 4, 6, 13, 16}
        {-, L, U, -, -, -, -, -, Acc, Size, 0, PPN, 0}
    \end{flushleft}
    
    The L bit will lock the entry in the TLB. When looking for a free entry, these entries are skipped. If all entries are locked entry 0, TLB entry zero is selected for replacement. The U bit will mark the address range described by the TLB entry as uncached.
    
    The access rights field is a four bit field that contains the page type and the privilege levels as shown in the access rights chapter. 
    
     The size field encodes the page size. It is 4-bit field encoding the supported page size: 
     	\[ \texttt{( 1U << ( pSize * 2 )) * T64\_PAGE\_SIZE} \] 
    A page size is not allowed to cross region boundaries and is aligned with selected the page size.
    Depending on the page size specified, the VPN and PPN fields are shorter and the offset is larger. The sum of both fields is however always 52 bits for virtual pages, and 40 bits for physical pages.
\end{iPageDescEntry}

\begin{iPageDescOpEntry}{Operation}
\begin{lstlisting}[style=iPageOpStyle]
insertIntoTlb( RegB, RegA );
\end{lstlisting}
\end{iPageDescOpEntry}

\begin{iPageDescEntry}{Exceptions}
    Privilege Violation Trap.
\end{iPageDescEntry}

\begin{iPageDescEntry}{Notes}
    None.  
\end{iPageDescEntry}
